%%%%%%%%%%%%%%%%%%%%%%%%%%%%%%%%%%%%%%%%%%%%%%%%%
\section{总结}
\label{sec:summary}

我们证明了准 PDF 的紫外发散行为与 PDF 的大不相同。准 PDF 的重正化是一个简单的乘法因子，而 PDF 的重正化是卷积形式，其原因在于 PDF 的紫外发散并非完全局域，因而不同味的 PDF 在重正化下相互混合。

我们证明了坐标空间准 PDF 微扰紫外发散的时空局域性，使得坐标空间准 PDF 在 QCD 微扰论的所有阶上都可以乘法重正化，如式~(\ref{eq:renor}) 所示。结合文献 \cite{Ma:2014jla} 中关于把准 PDF 全部领头幂次的共线（CO）发散因子化为 PDF 的全阶论证，我们的结论是：重正化后的准 PDF 是从 LQCD 计算提取 PDF 的良好候选者。

然而，对准 PDF 的 LQCD 计算而言，引入带微扰计算的 $C_i$ 的整体因子 $e^{-C_i |\xi_z|}$ 并不足以消除所有幂次发散，因为我们无法把 $C_j$ 计算到所有阶。也就是说，我们需要非微扰地重正化准 PDF 的幂次紫外发散。如文献 \cite{Ishikawa:2016znu} 所示，例如可以引入一个整体非微扰因子，在所有阶消除这类幂次发散，此处不展开细节。原则上，非微扰重正化准 PDF 幂次发散的方案并不唯一。但是，只要重正化是乘法的，重正化程序就不改变准 PDF 的共线（CO）性质（从而准 PDF 仍可因子化为 PDF）；不同的重正化程序／选择对应于不同的重正化方案，并将导致准 PDF 与 PDF 之间不同的匹配系数。

此外，我们证明了坐标空间准 PDF 对一切 $\xi_z$ 取值都表现良好，唯独 $\xi_z=0$ 除外。也就是说，如果想通过傅里叶变换从坐标空间准 PDF 得到动量空间准 PDF，就必须在变换之前定义一套一致的减除方案，以消除 $\xi_z\to 0$ 处的所有发散项。没有这一减除，动量空间准 PDF 在大 $\tilde{x}$ 区域是不适定的。

{\it 补充说明：}最近出现了两项关于准 PDF 重正化的独立研究 \cite{Ji:2017oey,Green:2017xeu}，虽然其方法与我们的很不一样，却得到了相同的结论。
